%% sections/03-methodology.tex

\section{Methodology}
\label{sec:methodology}

\subsection{The Marathon Workshop Pipeline}

The Marathon D\&D Workshop is a structured pipeline for designing and playtesting D\&D
5e adventures set in the Dragonlance campaign setting. Every session follows an
identical seven-stage sequence that produces a reproducible, fully logged session
artifact from a compiled adventure module and a character-sheet-specified party. The
stages are:

\begin{enumerate}
  \item \textbf{PREP.} The adventure is compiled to a module document that inlines all
  room descriptions, encounter stat blocks, treasure manifests, and critical rule text.
  The party's character sheets are loaded. The applicable rubric version is recorded in
  the session frontmatter. Dice seeds are initialized.

  \item \textbf{PLAY.} An AI system (Claude, Anthropic) runs the session as DM while
  simultaneously playing all player characters according to their Playstyle Heuristics ---
  per-PC behavioral specifications that govern when each character speaks, acts
  tactically, pursues personal goals, and defers to other party members. All dice rolls
  use a seeded pseudorandom number generator for reproducibility.

  \item \textbf{LOG.} The session log is produced as a structured narrative document
  recording all scenes, dice rolls with natural values and applied modifiers, PC actions
  and narration, and tagged SURPRISE moments where outcomes exceeded design prediction.

  \item \textbf{GATE.} The session log is scored against the locked rubric version
  using the session-gate protocol described in Section~\ref{sec:scoring-protocol},
  producing a per-dimension scorecard and an overall gate score out of 80.

  \item \textbf{PANEL.} Five player-experience lenses (Tactician, Roleplayer,
  Lorekeeper, Improviser, Fresh Face) each review the session log from their distinct
  perspective and produce a weighted aggregate panel score. Gate and panel scoring
  are performed independently.

  \item \textbf{INNOVATE.} The session log and panel reviews are scanned for outcomes
  that existing rubric anchors cannot adequately score. Qualifying outcomes are logged
  as innovations with a dimension tag. When two or more logged innovations cluster in
  the same dimension, an amendment proposal is drafted per the protocol in
  Section~\ref{sec:amendment-pipeline}.

  \item \textbf{HANDOFF.} A canonical handoff document is produced that records all
  campaign-permanent facts established in the session, updates the TRACKER with session
  metrics, and registers any new seeds or retrieved seeds in the campaign seed log.
  No session is considered complete without a handoff.
\end{enumerate}

\subsection{The 8-Dimension Rubric}
\label{sec:rubric-dimensions}

The rubric scores each session on eight dimensions, each scored 0--10 using four band
levels: 0--3 (significant failure), 4--6 (adequate but limited), 7--9 (strong
evidence of presence), and 10 (full anchor requirements met). The maximum session score
is 80. Dimensions are designed to be independent: a session's score on one dimension
should not mechanically constrain its score on another, except where explicit
cross-dimensional anchors have been ratified (see MF below).

\paragraph{Engagement (Eng).}
\textit{Did each player character contribute meaningfully to the session's outcomes?}
A score of 0--3 indicates that one or more PCs were passive observers for most of the
session, with no scene they initiated or driven by their character-specific motivation.
A score of 4--6 indicates that all PCs were present and acted but that their
contributions were primarily reactive rather than character-motivated. A score of 7--9
requires that every PC drove at least one scene and that each PC's contribution is
traceable to something specific on their character sheet. A score of 10 requires the
additional condition that the PCs' combined contributions produced an outcome that the
adventure's module did not pre-script --- a genuinely emergent collective result.

\paragraph{Mechanical Fairness (MF).}
\textit{Did dice rolls produce outcomes that felt honestly earned, and did mechanics
serve the fiction?}
The 0--3 band covers sessions where dice results were narratively arbitrary (high roll
produces low-value outcome, or vice versa) or where a mechanic functioned in a way
that undermined player choice. The 4--6 band covers sessions where dice and outcomes
were consistent but mechanics were applied narrowly, without narrative integration.
The 7--9 band requires that every significant dice result is integrated into the
narrative without retconning and that at least one mechanic is used in a way that
extends a character's established fictional positioning (e.g., an Athletics roll that
both succeeds mechanically and advances a character's established physical competence
in-fiction). The 10-band requires, in addition to the 7--9 criteria, a
\textit{cross-character mechanical compound}: at least one instance where one PC's
class ability or spell is applied to another PC's roll, framed as an in-character
action rather than metagame resource management. This anchor was added at v1.4 after
Session 6 of Campaign 1 produced the first instance of this phenomenon (a Divination
Wizard invoking her Portent die on behalf of a named ally while speaking in her
mentor's voice).

A separate MF anchor, ratified at v1.1, caps MF at 6 when any mechanic in the session
systematically damages outcomes in other dimensions. This \textit{cross-dimensional-harmlessness}
requirement reflects the finding from Session 1 that a Constitution-save exhaustion
mechanic caused two atmospheric payoffs to fail through mechanical side-effect rather
than narrative logic, producing false Atmospheric Landing failures that were not
captured by any initial anchor.

\paragraph{Pacing (Pac).}
\textit{Did scenes earn their length, and did the session's arc feel appropriately
shaped?}
The 0--3 band covers sessions where one or more scenes ran significantly past or short
of their narrative content --- a combat encounter that extended for three rounds after
the decisive action was resolved, or a revelation scene cut so short that its
emotional content had no time to register. The 4--6 band covers sessions with adequate
pacing in most scenes but one notable structural mismatch. The 7--9 band requires that
every scene's length is proportional to its narrative weight, including transition
scenes and downtime scenes. The 10-band adds the requirement that the session's
overall arc is visible in the session log: a reader with no prior knowledge of the
campaign should be able to identify a beginning, a middle, and an end to the session
as a self-contained narrative unit.

\paragraph{Character Agency (CA).}
\textit{Did player characters make choices that the module required them to make, and
did those choices shape the outcome?}
The 0--3 band covers sessions where outcomes were determined by module logic regardless
of what PCs chose. The 4--6 band covers sessions where PCs chose from a menu of
options but the options were all equivalent in narrative consequence. The 7--9 band
requires that at least one session outcome depended specifically on a PC choice that
the module anticipated but did not require --- the module offered the choice; the PC
took it; and taking it led to a materially different outcome than not taking it would
have produced. The 10-band was amended twice (v1.9, v1.10) to add two additional
qualifying paths: a \textit{sustained behavioral pattern} (a sequence of non-actions
across multiple scenes that collectively constitutes the choice, with each individual
non-action ambiguous but the pattern deliberate) and a \textit{character-register
conversion} (a PC converts a mechanical situation into a character-authored outcome
traceable to their character sheet heuristics rather than tactical calculation, such as
ending a combat encounter by naming an NPC's coercion rather than dealing damage).

\paragraph{Module Fidelity (MFid).}
\textit{Did the adventure deliver the outcomes and payoffs it was designed to deliver?}
The 0--3 band covers sessions where one or more designed payoffs failed entirely ---
the artifact was not found, the NPC arc did not fire, the climactic scene was skipped.
The 4--6 band covers sessions where designed payoffs occurred but in reduced form.
The 7--9 band requires that all designed anchor symptoms were present in the session.
The 10-band adds the requirement that at least one hidden payoff --- an outcome the
module designed but did not make explicit in the anchor list --- was also delivered.
Amendment v1.2 added a \textit{passive-Perception fallback}: if a primary anchor symptom
fails because no PC rolled high enough to trigger it, a passive-Perception delivery of
the same information (via environmental description rather than active investigation)
counts as a valid anchor-symptom delivery. Amendment v1.3 added a
\textit{positional-redundancy} requirement: at the 7--9 band and above, anchor symptoms
must have two of three fallback delivery paths available in the session log, ensuring
that module design provides sufficient redundancy rather than placing all delivery
weight on a single roll.

\paragraph{Atmospheric Landing (AL).}
\textit{Did designed emotional and sensory moments land as affecting?}
Atmospheric Landing received the most amendments of any dimension (v1.1, v1.2, v1.4,
v1.5, v1.6, v1.7, v1.8, v1.9, v1.12) because novel positive outcomes most frequently
exceeded its initial anchors. The 0--3 band covers sessions where designed atmospheric
moments are present in the log but are not integrated into PC responses --- the
description is there but no PC acknowledgment, silence, or behavioral change follows.
The 4--6 band covers sessions where at least one atmospheric moment is acknowledged
by a PC but the acknowledgment is brief and does not persist into subsequent scenes.
The 7--9 band requires one moment that shows evidence of genuine impact: either a
specific sensory or emotional image that is subsequently referenced (verbally or
behaviorally) by at least one PC, or a moment of sustained silence or non-response
that the session log itself flags as significant.

The 10-band evolved through successive amendments to encompass six distinct qualifying
paths, any one of which suffices: (a) a specific image that will outlive the session,
acknowledged by at least one PC through an inter-PC chain (v1.1); (b) a multi-scene
silent-reception chain where a PC holds an unarticulated understanding across three or
more scenes and acts on it without naming it (v1.4); (c) an act-without-announcement
where a PC acts differently as a direct consequence of a prior naming event, without
declaring the connection (v1.5); (d) simultaneous arc-convergence where a single naming
act triggers multiple outstanding arc-completions in distinct PCs (v1.6); (e) a
framework-failure arc-activation where a supporting NPC's expected role collapses and a
PC steps into the resulting gap without prompting (v1.7); and (f) an intra-session
restraint anchor where a PC demonstrates the same act-without-announcement behavior
within the same session (rather than cross-session), signaling a behavioral change
still in the moment of formation (v1.12).

\paragraph{Surprise (Sur).}
\textit{Did anything genuinely emergent happen that the adventure design could not have
predicted?}
The Surprise dimension explicitly rewards outcomes that were not designed. The 0--3
band covers sessions with no logged surprise moments --- every outcome was fully
predictable from the module and the character sheets. The 4--6 band covers sessions
with one or two minor surprises (unexpected dice results, minor improvised NPC
responses). The 7--9 band requires at least one substantive emergent outcome that
changed the adventure's trajectory or produced a narrative development the module
did not anticipate. The 10-band requires that this outcome is not only substantive
but amendment-worthy: it represents a quality that existing rubric anchors cannot
capture, or it demonstrates that a party archetype produces session dynamics not
previously seen in the corpus.

\paragraph{Table Readiness (TR).}
\textit{Could a DM pick up this session's compiled module and run it without external
references?}
This dimension evaluates the module and session infrastructure rather than the session
narrative. The 0--3 band covers sessions where key rule text was absent from the
compiled module (requiring lookup during play), NPCs lacked stat blocks, or encounter
balance was unvalidated. The 4--6 band covers sessions where critical references were
inline but supplementary rules required lookup. The 7--9 band requires that all rule
text necessary to run every encounter, trap, and hazard in the session is inlined in
the module document, NPCs have complete stat blocks, and encounter balance has been
validated against party tier and size. The 10-band adds the requirement that the
module's DM cheatsheet includes the three most likely player-decision branches and
their adjudication, enabling the DM to handle the session's principal choice points
without improvisation under pressure.

\subsection{The Amendment Pipeline}
\label{sec:amendment-pipeline}

The rubric's amendment mechanism is the instrument through which empirical use refines
the quality space it can measure. The pipeline has five steps:

\begin{enumerate}
  \item \textbf{Innovation identification.} At the INNOVATE stage of the session
  pipeline, the session log and panel reviews are read with the question: did any
  session outcome exceed the predictive power of the current rubric anchors? An outcome
  that is clearly captured (positive or negative) by existing anchor language does not
  qualify. An outcome that required the reviewer to reason outside the rubric's anchor
  set qualifies as an innovation candidate.

  \item \textbf{Innovation logging.} Qualifying outcomes are logged to the innovations
  register with a dimension tag (identifying which rubric dimension is most relevant),
  a scope classification (anchor-gap: the dimension's existing anchors cannot capture
  the outcome; or new-dimension-candidate: the outcome suggests a quality orthogonal to
  all eight existing dimensions), and a verbatim log citation.

  \item \textbf{Clustering and proposal.} When two or more logged innovations carry the
  same dimension tag, an amendment proposal is drafted. The proposal describes the
  pattern in generalized terms (not specific to the source session), specifies the band
  to which the new anchor applies, and justifies why the pattern represents a genuinely
  new form of the dimension's quality rather than a restatement of existing anchors.

  \item \textbf{Ratification.} The proposal is reviewed against three criteria: the
  pattern is genuinely novel (not already captured by existing anchor language through
  reasonable interpretation); the pattern is reproducible (it could in principle appear
  in future sessions with different parties and adventures); and the proposed anchor
  language is precise enough to enable consistent scoring. If all three criteria are
  met, the proposal is adopted, the rubric version is incremented, and the constituent
  innovations are marked as adopted in the register.

  \item \textbf{Forward-only application.} The new anchor applies to all sessions
  scored after the amendment date. Sessions scored under prior versions are not
  retroactively rescored. Cross-version comparisons are made explicit by reporting the
  rubric version alongside gate scores.
\end{enumerate}

This forward-only constraint preserves the integrity of the historical score series:
each session score reflects the quality space that was operationally defined at the
time it was played. Retrospective revision would conflate the rubric's learning with
genuine quality change, making it impossible to distinguish a 12-point improvement
in session quality from a 12-point expansion of the rubric's scoring ceiling.

\subsection{Session-Gate Scoring Protocol}
\label{sec:scoring-protocol}

For each dimension, the session-gate reviewer reads the session log with that
dimension as sole focus. The reviewer identifies supporting evidence (specific quoted
passages from the log) and counter-evidence (absences, failures, or degraded outcomes),
assigns a score using the band descriptions, and writes an evidence cell that documents
the specific log citations on which the score rests. The scoring protocol explicitly
prohibits generous or harsh adjustment: ``Do not score generously out of sympathy. Do
not score harshly to seem rigorous. Score what is in the log.''

The binding pass threshold is 56/80 (70\%), set conservatively to be achievable by a
well-designed session scoring 7/10 on all eight dimensions without requiring
near-perfect performance. Sessions 1--3 of each campaign are scored as advisory only,
because adventures during this window are still under active revision and the session
may be the first time the module has been run. Advisory scores contribute to the
innovation log but do not determine pass/fail. Sessions 4 onward are binding.

\subsection{The Seven Campaigns}

The corpus consists of seven complete campaigns, each comprising seven adventures and
seven sessions, for a total of 49 sessions. Campaigns were designed with distinct
party archetypes and central dramatic questions to test whether the rubric's scores
depend on thematic content or reflect session execution quality independent of theme.

Campaign 1 (Moon-Silver Cycle) used a grief-themed party with explicit loss-hooks on
every character sheet, centered on a campaign question of what an emotion is worth
trading away. Campaign 2 (Conclave Compact) used a covenant-restoration party
structure with a convergence deadline generating mechanical time pressure. Campaign 3
(Halted Spire) used a faith-and-authority theme with a party defined by a crisis of
institutional legitimacy. Campaign 4 (Thorngate Watch) introduced resource exhaustion
mechanics (Supply Die, no long rests for seven sessions) with a siege-and-resource
management archetype. Campaign 5 (Thieves Guild) used a professional-code motivation
with no grief hooks and no designed personal stakes, testing whether the rubric's
Atmospheric Landing dimension depends on grief-adjacent party design. Campaign 6
(Military Unit) used a duty-spine with an institutional deliverable. Campaign 7 (Party
of Strangers) used a party with no designed personal stakes, no grief paragraphs, no
loyalty objects, and no professional codes --- the most stripped-down archetype in the
corpus, designed to test whether the rubric can score sessions when parties have no
designed emotional infrastructure.
